\chapter{智慧水利三维场景构建}

\paragraph*{学习目标}

通过本章学习，学生应能够：
\begin{enumerate}
\item 说明三维渲染坐标变换与WebGL2渲染管线，使用Three.js加载工程模型；
\item 区分WMS、WMTS与WFS，正确处理CGCS2000、投影坐标和高程基准；
\item 说明倾斜摄影、BIM与3D Tiles在水利场景中的分工和生产流程；
\item 设计数字孪生水利平台的分层架构、数据闭环和验收指标；
\item 将三维场景与实时监测、模型计算和业务服务建立可追踪关联。
\end{enumerate}

\paragraph*{引言}

三维场景不是智慧水利平台的装饰层，而是空间基准、工程模型、监测数据和业务状态的综合表达。本章按照“浏览器三维渲染—GIS与坐标—实景/BIM模型—数字孪生平台”的顺序展开，技术基线为WebGL2、Three.js r160+、Cesium 1.11x、OGC服务与IFC 4.3。

\section{WebGL2与Three.js三维渲染基础}

\subsection{坐标变换与透视投影}

一个顶点通常依次经过模型坐标系、世界坐标系、观察坐标系、裁剪坐标系、规范化设备坐标和屏幕坐标。模型矩阵$M$把局部构件放入世界，观察矩阵$V$表示相机，投影矩阵$P$形成透视或正交投影，三者合称MVP矩阵：
\[
  \mathbf{p}_{clip}=PVM\mathbf{p}_{model}.
\]

设垂直视场角为$\theta$、宽高比为$a$、近远裁剪面距离为$n$和$f$，OpenGL风格透视矩阵可写为
\[
P=\begin{bmatrix}
\dfrac{1}{a\tan(\theta/2)}&0&0&0\\
0&\dfrac{1}{\tan(\theta/2)}&0&0\\
0&0&-\dfrac{f+n}{f-n}&-\dfrac{2fn}{f-n}\\
0&0&-1&0
\end{bmatrix}.
\]
这里$f$只表示远裁剪面，不能同时当作焦距。近裁剪面过小会降低深度缓冲精度；大尺度库区场景应按可见范围动态调整近远面，并采用局部原点或高低位拆分缓解浮点精度问题。

\subsection{WebGL2渲染管线}

本章统一使用WebGL2，不提供名义上的WebGL1回退。WebGL1对顶点数组对象、32位索引和扩展支持的要求不同，若确需兼容，应单独实现并测试兼容分支。

\begin{lstlisting}[language=JavaScript,breaklines=true]
const canvas = document.querySelector('#scene');
const gl = canvas.getContext('webgl2');
if (!gl) throw new Error('当前浏览器不支持 WebGL2');

gl.enable(gl.DEPTH_TEST);
gl.clearColor(0.86, 0.93, 0.98, 1.0);
gl.viewport(0, 0, canvas.width, canvas.height);
gl.clear(gl.COLOR_BUFFER_BIT | gl.DEPTH_BUFFER_BIT);
\end{lstlisting}

顶点着色器必须显式声明使用的uniform。下面的最小着色器把模型、观察和投影矩阵分开，便于定位坐标错误。

\begin{lstlisting}
#version 300 es
layout(location = 0) in vec3 aPosition;
uniform mat4 uModel;
uniform mat4 uView;
uniform mat4 uProjection;

void main() {
    gl_Position = uProjection * uView * uModel
                  * vec4(aPosition, 1.0);
}
\end{lstlisting}

CPU负责准备缓冲、纹理和绘制命令，GPU执行顶点处理、图元装配、光栅化、片元处理和深度/混合测试。性能差异受场景、设备和实现影响，不使用“固定快10倍或100倍”之类无出处断言；应通过帧时间、绘制调用数和显存占用实测。

\subsection{Three.js场景组织与模型加载}

Three.js封装了场景、相机、材质、灯光和渲染器。工程项目宜把测站、坝段、闸门等对象组织为稳定层级，并把业务标识写入\texttt{userData}，而不是依赖模型节点的显示名称。

官方addon路径为\texttt{three/addons/loaders/GLTFLoader.js}\cite{threeGLTFLoader}。glTF/GLB适合Web交付；加载后仍需检查单位、坐标轴、材质色彩空间和包围盒。

\begin{lstlisting}[language=JavaScript,breaklines=true]
import * as THREE from 'three';
import { GLTFLoader } from
  'three/addons/loaders/GLTFLoader.js';

const scene = new THREE.Scene();
const loader = new GLTFLoader();
const gltf = await loader.loadAsync('/models/dam.glb');
gltf.scene.userData.assetId = 'DAM-001';
scene.add(gltf.scene);
\end{lstlisting}

模型交付前应完成网格简化、纹理压缩、法线检查和层级拆分。优化顺序是先测量瓶颈，再选择实例化、LOD、视锥裁剪、遮挡裁剪或按需加载；不同策略的效果必须以目标终端实测为准。

\section{GIS服务、坐标与二三维集成}

\subsection{OGC地图与要素服务}

WMS根据请求范围、尺寸、坐标参考系和样式动态渲染地图图像；WMTS按预定义瓦片矩阵提供地图瓦片，便于缓存和大规模浏览；WFS提供可查询的矢量要素。WMS不是“栅格瓦片服务”，WMTS才以预定义瓦片为核心\cite{ogcWMS,ogcWMTS}。

\begin{table}[htbp]
\centering
\caption{常用OGC服务的职责}
\label{tab:ogc-services}
\begin{tabular}{llll}
\toprule
服务 & 返回内容 & 主要特点 & 适用场景 \\
\midrule
WMS & 动态地图图像 & 范围与样式灵活 & 专题图叠加 \\
WMTS & 预定义地图瓦片 & 易缓存、吞吐高 & 稳定底图 \\
WFS & 矢量要素 & 可查询属性与几何 & 河道、测站编辑查询 \\
\bottomrule
\end{tabular}
\end{table}

教学示例使用自建GeoServer，不依赖商业底图密钥：

\begin{lstlisting}[language=JavaScript]
const imagery = new Cesium.WebMapTileServiceImageryProvider({
  url: 'https://gis.example.edu/geoserver/gwc/service/wmts',
  layer: 'water:cgcs2000_basemap',
  style: '',
  format: 'image/png',
  tileMatrixSetID: 'EPSG:4490'
});
viewer.imageryLayers.addImageryProvider(imagery);
\end{lstlisting}

实际部署需要将示例域名替换为课程或项目服务器，并在服务端配置跨域、访问控制和缓存策略。

\subsection{CGCS2000与投影坐标}

CGCS2000地理二维坐标系代码是EPSG:4490\cite{epsg4490}，EPSG:4326对应WGS 84。CGCS2000参考椭球长半轴为6378137米、扁率倒数为298.257222101；它与WGS 84椭球参数极为接近，但不能在坐标定义中直接写成\texttt{datum=WGS84}。

我国6度高斯—克吕格分带约为13至23带。经度$\lambda$所在带号可按
\[
  N=\left\lfloor\frac{\lambda}{6}\right\rfloor+1,
  \qquad \lambda_0=6N-3
\]
确定中央经线。黄河、长江等流域跨越多个带，应依据具体数据经度选择分带或使用项目统一投影，不能硬编码“黄河19–21带、长江28–30带”。

下面示例把CGCS2000经纬度转换到中央经线111度的6度带（19带）。\texttt{tmerc}表示横轴墨卡托/高斯—克吕格计算；带号前缀是否写入东坐标，应与项目数据规范一致。

\begin{lstlisting}[language=JavaScript]
const CGCS2000 =
  '+proj=longlat +a=6378137 +rf=298.257222101 +type=crs';
const GK19 =
  '+proj=tmerc +lat_0=0 +lon_0=111 +k=1 '
  + '+x_0=19500000 +y_0=0 +a=6378137 '
  + '+rf=298.257222101 +units=m +type=crs';

proj4.defs('CGCS2000', CGCS2000);
proj4.defs('CGCS2000_GK19', GK19);
const projected = proj4('CGCS2000', 'CGCS2000_GK19',
                        [110.50, 34.20]);
\end{lstlisting}

\subsection{高程基准与空间一致性}

GNSS通常给出椭球高，水利工程常使用1985国家高程基准下的正常高。二者通过高程异常联系，术语不是“加速度重力异常”。在我国，高程异常可达到十几米至几十米量级；淹没线、闸顶和监测点高程必须记录基准、单位、转换模型和版本，不能只保存一个数值。

二三维集成应建立统一空间元数据：水平CRS、高程基准、时间基准、单位、精度等级和转换记录。对关键坝轴线、控制点和监测点进行同名点校核，发现系统偏差时回到数据生产环节处理，而不是在前端用手工偏移掩盖。

\subsection{Cesium场景集成}

Cesium的\texttt{Cesium3DTileset.fromUrl}返回Promise，应等待对象创建完成后再加入场景\cite{cesium3DTileset}：

\begin{lstlisting}[language=JavaScript]
try {
  const tileset = await Cesium.Cesium3DTileset.fromUrl(
    '/tiles/dam/tileset.json'
  );
  viewer.scene.primitives.add(tileset);
  await viewer.zoomTo(tileset);
} catch (error) {
  console.error('三维瓦片加载失败', error);
}
\end{lstlisting}

业务图层采用稳定标识关联测站与工程对象。点击三维构件后，前端用标识查询后端状态；数据更新只改变颜色、标签或图表，不重复创建整个模型。

\section{倾斜摄影与BIM模型构建}

\subsection{倾斜摄影数据生产}

倾斜摄影通过垂直与多个倾斜视角获取影像，用于恢复工程及周边地表。系统由飞行平台、航摄仪、地面控制和内业处理组成，如图\ref{fig:uav-oblique-system}所示。

\begin{figure}[htbp]
\centering
\uavObliquePhotographySystem
\caption{无人机倾斜摄影测量系统组成}
\label{fig:uav-oblique-system}
\end{figure}

生产流程包括航线设计、影像采集、质量检查、空中三角测量、密集匹配、网格与纹理生成、坐标校核、切片和发布。软件名称统一使用Agisoft Metashape与Bentley ContextCapture；不再混用PhotoScan、Smart 3D等旧称。

\begin{figure}[htbp]
\centering
\obliqueImagePreprocessQC
\caption{倾斜影像预处理与质量控制流程}
\label{fig:oblique-preprocess}
\end{figure}

术语必须区分：特征点由算法在影像中检测；连接点是多幅影像中匹配到同一地物的特征；像控点具有外业测得的已知坐标，用于约束空三结果；检查点不参与平差，只用于独立精度评定。

若检查点共有$n$个，成果与实测坐标差分别为$\Delta X_i$、$\Delta Y_i$、$\Delta H_i$，三个方向的均方根误差应分别计算：
\begin{equation}
  m_x=\sqrt{\frac{1}{n}\sum_{i=1}^{n}(\Delta X_i)^2}.
  \label{eq:rmse-x}
\end{equation}
\begin{equation}
  m_y=\sqrt{\frac{1}{n}\sum_{i=1}^{n}(\Delta Y_i)^2}.
  \label{eq:rmse-y}
\end{equation}
\begin{equation}
  m_h=\sqrt{\frac{1}{n}\sum_{i=1}^{n}(\Delta H_i)^2}.
  \label{eq:rmse-h}
\end{equation}
精度结论应同时给出检查点数量、空间分布、坐标/高程基准和限差来源，不能只报告一个平均值。

\subsection{BIM语义与IFC 4.3}

BIM强调构件语义、属性和工程关系，倾斜摄影强调表面真实感，两者不能互相替代。闸门启闭机、监测仪器等关键对象宜保留可查询语义；大范围地形和周边建筑可用倾斜摄影或地形瓦片表达。

IFC 4.3新增Ports and Waterways Domain，覆盖水道、船闸、闸门、溢洪道等部分设施语义\cite{ifc43Ports}。不能用\texttt{IfcDoor}表示闸门、\texttt{IfcFurniture}表示廊道或把坝体简单当作\texttt{IfcStructuralMember}。同时必须说明标准边界：buildingSMART当前文档仍把坝、堤、堰列为该域范围之外。对这些对象应采用项目分类、属性集和映射规则，并保留原始设计模型，不虚构不存在的IFC实体。

模型轻量化应保留业务所需的全局标识、类型、空间层级和关键属性，再把几何转换为glTF或3D Tiles。转换日志记录源模型版本、转换工具、参数、时间和校核结果，以便发现属性丢失时追溯。

\subsection{多源模型融合与发布}

融合前先统一坐标、高程、单位和时间版本，再处理几何。推荐步骤为：
\begin{enumerate}
\item 建立GIS空间底座与项目局部坐标转换；
\item 校核倾斜摄影成果的控制点和边界；
\item 对BIM构件建立业务编码与空间定位；
\item 生成分层LOD和3D Tiles/glTF交付物；
\item 通过同名点、剖面和关键尺寸进行复核；
\item 发布模型目录、元数据与版本清单。
\end{enumerate}

模型成果分析包括几何精度、纹理质量、语义完整性、加载性能和版本一致性。任何优化都不应破坏测站—构件—业务记录之间的关联。

\section{数字孪生水利平台架构}

\subsection{概念边界与五维组织}

数字模型是物理对象的数字表达；数字影子强调物理对象向数字模型的数据更新；数字孪生还要求虚实之间形成受控、可验证的双向业务闭环。仅有精美三维场景而没有持续数据、模型服务和反馈流程，不能称为完整数字孪生。

NASA文献强调物理模型、传感器更新和运行历史的综合映射\cite{glaessgen2012}。陶飞、张萌的数字孪生车间模型包含物理空间、虚拟空间、服务系统和孪生数据\cite{tao2017dts}。为便于水利平台设计，本教材把连接机制单独列为第五维，形成
\[
  DT=(PE,VE,Ss,DD,CN),
\]
其中$PE$为物理实体、$VE$为虚拟实体、$Ss$为服务、$DD$为孪生数据、$CN$为连接。该五维写法是对原始构成的工程化展开，不把“孪生数据”改写成“数据处理”。

\begin{figure}[htbp]
\centering
\begin{tikzpicture}[>=Stealth,node distance=1.2cm and 1.4cm]
\tikzset{d/.style={draw=Primary,rounded corners=2pt,fill=PrimaryLight,
  minimum width=3.0cm,minimum height=.8cm,align=center,text=PrimaryDark}}
\node[d] (pe) {物理实体 $PE$\\工程、设备、环境};
\node[d,right=of pe] (ve) {虚拟实体 $VE$\\几何、机理、状态};
\node[d,below=of pe] (dd) {孪生数据 $DD$\\实时、历史、元数据};
\node[d,right=of dd] (ss) {服务 $Ss$\\分析、预警、调度};
\node[d,below right=.45cm and -.15cm of dd] (cn) {连接 $CN$\\感知、接口、反馈};
\draw[<->,draw=Primary,thick] (pe)--(ve);
\draw[<->,draw=Primary,thick] (pe)--(dd);
\draw[<->,draw=Primary,thick] (ve)--(ss);
\draw[<->,draw=Primary,thick] (dd)--(ss);
\draw[->,draw=Accent,thick] (cn)--(pe);
\draw[->,draw=Accent,thick] (cn)--(ve);
\draw[->,draw=Accent,thick] (cn)--(dd);
\draw[->,draw=Accent,thick] (cn)--(ss);
\end{tikzpicture}
\caption{数字孪生水利平台五维组织}
\label{fig:dt-five-dimensions}
\end{figure}

\begin{table}[htbp]
\centering
\caption{五个维度在大坝安全监测中的映射}
\label{tab:dt-five-dimensions}
\begin{tabular}{p{2.2cm}p{4.2cm}p{6.2cm}}
\toprule
维度 & 典型内容 & 设计关注点 \\
\midrule
物理实体 & 坝体、廊道、传感器、闸门 & 对象标识、状态、可控边界 \\
虚拟实体 & 三维几何、渗流/变形模型、规则 & 适用条件、参数、版本、可信度 \\
服务 & 数据质检、预警、分析、方案比选 & 输入输出、时效、责任主体 \\
孪生数据 & 实时值、历史序列、模型结果、档案 & 时空基准、质量码、血缘 \\
连接 & 采集协议、API、消息、反馈指令 & 安全、延迟、幂等、审计 \\
\bottomrule
\end{tabular}
\end{table}

\subsection{分层技术架构}

平台可划分为物理对象层、感知连接层、数据底板层、模型与知识层、服务层和交互应用层。分层的目的不是增加名词，而是明确数据在哪里校验、模型在哪里运行、结果由谁解释、指令由谁批准。

\begin{figure}[htbp]
\centering
\begin{tikzpicture}[node distance=.35cm]
\tikzset{layer/.style={draw=Primary,rounded corners=2pt,fill=PrimaryLight,
  minimum width=12.2cm,minimum height=.75cm,align=center,text=PrimaryDark}}
\node[layer] (app) {交互应用层：三维态势、告警处置、方案比选、移动巡检};
\node[layer,below=of app] (service) {服务层：查询、预警、分析、推演、调度建议、权限与审计};
\node[layer,below=of service] (model) {模型与知识层：水文/水力/结构模型、规则库、模型注册与评估};
\node[layer,below=of model] (data) {数据底板层：时序、空间、工程档案、主数据、元数据与血缘};
\node[layer,below=of data] (connect) {感知连接层：网关、数据质检、消息总线、API、时钟同步};
\node[layer,below=of connect] (physical) {物理对象层：流域、河网、水库、水闸、传感器与执行设备};
\draw[<->,draw=Primary,thick] (app)--(service);
\draw[<->,draw=Primary,thick] (service)--(model);
\draw[<->,draw=Primary,thick] (model)--(data);
\draw[<->,draw=Primary,thick] (data)--(connect);
\draw[<->,draw=Primary,thick] (connect)--(physical);
\end{tikzpicture}
\caption{数字孪生水利平台分层架构}
\label{fig:dt-layered-architecture}
\end{figure}

感知连接层先完成测站身份、时间戳、单位、范围和质量码检查，再进入数据底板。模型层不直接读取任意原始表，而通过具有版本的数据产品获得输入。服务层把模型结果转换为可解释的风险等级或方案指标；应用层展示依据与不确定性，不把建议伪装成自动命令。

\subsection{业务场景与需求分解}

数字孪生建设应从闭环业务场景开始，而不是从“建立全域高精度模型”开始。场景说明至少回答六个问题：管理对象是什么、谁在何时使用、需要哪些观测、调用什么模型、输出如何解释、最终行动由谁批准。不同场景的时效和空间精度要求差异很大，不能共享一套未经区分的“实时”指标。

\begin{table}[htbp]
\centering
\caption{典型数字孪生水利场景的需求差异}
\label{tab:dt-scenario-requirements}
\begin{tabular}{p{2.6cm}p{3.4cm}p{3.5cm}p{3.5cm}}
\toprule
场景 & 核心对象与观测 & 主要模型/规则 & 业务输出 \\
\midrule
大坝安全 & 坝段、测点、库水位、温度 & 基线、统计或结构模型 & 风险证据与处置工单 \\
洪水预报 & 河网、断面、雨量与流量 & 水文/水动力模型 & 过程预报与影响范围 \\
水资源调度 & 水库、取用水户、控制断面 & 供需平衡与优化模型 & 可行方案及约束冲突 \\
河湖监管 & 河湖岸线、遥感影像、巡查事件 & 变化检测与规则 & 疑点清单和核查任务 \\
工程巡检 & 构件、缺陷、工单、现场影像 & 缺陷分类与维护规则 & 定位、等级和复检计划 \\
\bottomrule
\end{tabular}
\end{table}

需求分解可采用“对象—事件—状态—服务—证据”模板。例如“大坝渗压异常”对象是测点和坝段，事件是新读数到达，状态包括质量码和风险级别，服务包括质检、模型计算和告警，证据包括原始值、趋势、相关测点和模型版本。模板能防止三维界面与后端业务脱节。

每个场景还应定义最小闭环。若第一阶段只能完成感知、质检、展示和人工处置，就应明确称为第一阶段能力，而不是用未来规划补足当前缺失。范围清楚比一次性堆叠全部模型更有利于验收和迭代。

\subsection{数据底板与模型平台}

数据底板至少管理四类数据：实时监测序列、空间地理数据、工程结构与档案、模型输入输出。每条关键数据应带有对象标识、采集时间、入库时间、单位、空间参考、质量码和来源。

\begin{table}[htbp]
\centering
\caption{数字孪生数据与模型注册的关键元数据}
\label{tab:dt-metadata}
\begin{tabular}{p{3cm}p{4.2cm}p{5.6cm}}
\toprule
对象 & 必备元数据 & 作用 \\
\midrule
监测数据集 & 测点、单位、时区、质量码、血缘 & 防止错点、错时、错单位 \\
空间数据集 & 水平CRS、高程基准、精度、版本 & 保证图层和模型空间一致 \\
计算模型 & 版本、参数、适用范围、校准记录 & 判断模型能否用于当前场景 \\
模型运行 & 输入快照、代码版本、开始/结束时间 & 支持复现和责任追踪 \\
业务规则 & 阈值、审批人、生效期、依据 & 避免规则无来源或长期失效 \\
\bottomrule
\end{tabular}
\end{table}

模型平台负责注册、调度、监控和评估，不把所有算法塞入同一服务。机理模型、统计模型和规则可以并存，但每个输出必须保留来源。模型更新后先回放历史工况并通过验收，再替换生产版本。

\subsection{模型卡与可信度管理}

模型“能够运行”不等于“适合决策”。每个生产模型应配模型卡，记录用途、责任人、输入、输出、适用空间与工况、参数来源、校准数据、误差指标、已知限制、版本和回滚方法。统计或机器学习模型还需记录训练数据时间范围和分布漂移检查。

\begin{table}[htbp]
\centering
\caption{水利模型卡的最小字段}
\label{tab:model-card}
\begin{tabular}{p{3.0cm}p{4.2cm}p{5.2cm}}
\toprule
字段组 & 示例 & 验证问题 \\
\midrule
身份与责任 & 模型ID、版本、负责人 & 谁批准进入生产环境 \\
适用范围 & 流域、工程、洪水量级、季节 & 当前工况是否超出范围 \\
输入契约 & 数据集版本、单位、时间步长 & 缺测或延迟如何处理 \\
输出与不确定性 & 流量过程、置信区间、风险等级 & 用户能否理解误差边界 \\
校准与验证 & 历史事件、指标、阈值 & 是否使用独立验证数据 \\
运行与回滚 & 资源、超时、降级、旧版本 & 失败后如何恢复服务 \\
\bottomrule
\end{tabular}
\end{table}

模型评估不宜只看一个平均误差。洪峰预报还可检查峰值、峰现时间和过程拟合；异常识别需结合误报、漏报和处置代价；空间淹没结果需检查边界位置和高程敏感性。所有指标都应与业务用途对应。

模型运行时先校验输入契约。若数据时间落后、单位不符或关键测站缺测，平台应拒绝运行、使用经批准的降级输入或显式降低可信等级，不能悄悄以零值填充。输出页面同时展示模型版本、输入时刻和可信状态，使使用者知道结果的条件。

模型变更遵循“开发—离线验证—影子运行—评审—发布—监控—回滚”流程。影子运行让新模型接收真实输入但不影响生产决策，可比较新旧结果并发现边界工况。模型审批记录与软件发布记录共同构成审计证据。

\subsection{事件时间、状态估计与质量码}

多源监测数据不会严格同时到达。平台应区分事件时间与处理时间：事件时间表示现场观测发生的时刻，处理时间表示平台收到或计算的时刻。模型输入按事件时间对齐；处理时间用于衡量传输和计算延迟。若只保留一个时间戳，补传数据可能被误认为当前状态。

流式处理可使用水位线（watermark）表达“系统认为某一事件时间之前的数据基本到齐”。水位线不是业务水位，而是事件时间处理术语。允许迟到窗口应依据网络和场景设置：秒级告警与日尺度统计不能共用同一窗口。超过窗口的数据仍应保存，但需触发重算、修订或标记为历史补录。

不同频率数据进入模型前需采用明确的对齐规则。连续物理量可在允许间隔内插值，离散设备状态通常保持最近有效值，累计雨量需要按时间段重采样。任何填补都要保留“原始、插值、估算或人工修订”标记，避免派生值被当作实测值。

\begin{table}[htbp]
\centering
\caption{监测数据质量码与处理建议}
\label{tab:dt-quality-flags}
\begin{tabular}{p{2.4cm}p{4.5cm}p{5.8cm}}
\toprule
质量状态 & 判断示例 & 下游处理 \\
\midrule
有效 & 格式、范围、时序和设备状态均正常 & 进入状态估计和模型 \\
可疑 & 突变、邻点不一致或接近量程边界 & 保留并降低权重，提示复核 \\
缺测 & 期望窗口内无观测 & 按模型契约停算或使用批准的估算 \\
迟到 & 事件时间早于当前水位线 & 保存并决定是否重算历史状态 \\
人工修订 & 经授权人员更正 & 保存原值、修订值、原因与操作者 \\
设备故障 & 自检失败、离线或维护 & 禁止作为正常观测参与模型 \\
\bottomrule
\end{tabular}
\end{table}

状态估计不能覆盖原始记录。原始层保持不可变观测，标准层完成单位和时空统一，状态层形成面向业务的当前状态，模型层保存预测或分析结果。各层通过数据血缘关联，用户可从风险结论追溯到模型运行、输入快照和原始测点。

单位转换在数据标准层集中完成。例如水位统一为米、流量统一为立方米每秒；接口仍携带单位字段。阈值和模型参数必须声明所用单位，禁止依赖开发者记忆。时区统一存储为带偏移的时间或UTC，界面按用户时区显示。

\subsection{三维交互与决策证据链}

三维界面应帮助用户定位对象、理解状态和获取证据。对象选择后显示稳定业务名称、编码、数据时刻和质量状态；颜色表达必须配合图例、文字或图标，不能只用红绿色区分风险。对键盘操作、字号和对比度也应进行基本无障碍检查。

状态着色与几何材质分离。原始模型材质属于资产表现，风险着色属于临时业务覆盖层；退出专题模式后应恢复原材质。这样既避免反复修改模型文件，也能同时支持渗压、位移、巡检等多个专题。

\begin{figure}[htbp]
\centering
\begin{tikzpicture}[>=Stealth,node distance=.75cm and 1.0cm]
\tikzset{evidence/.style={draw=Primary,rounded corners=2pt,fill=PrimaryLight,
 minimum width=3.0cm,minimum height=.7cm,align=center,text=PrimaryDark}}
\node[evidence] (object) {三维对象\\坝段/测点};
\node[evidence,right=of object] (state) {当前状态\\值、时刻、质量码};
\node[evidence,right=of state] (trend) {历史与关联\\趋势、邻点、工况};
\node[evidence,below=of state] (model) {模型证据\\版本、输入、残差};
\node[evidence,right=of model] (work) {业务闭环\\告警、工单、复核};
\draw[->,draw=Primary,thick] (object)--(state);
\draw[->,draw=Primary,thick] (state)--(trend);
\draw[->,draw=Primary,thick] (state)--(model);
\draw[->,draw=Primary,thick] (trend)--(work);
\draw[->,draw=Primary,thick] (model)--(work);
\end{tikzpicture}
\caption{从三维对象到业务处置的证据链}
\label{fig:dt-evidence-chain}
\end{figure}

时间轴允许回看任一历史状态，但必须区分“当时已知的数据”和“后来补录后重算的状态”。若回放使用了后来数据，界面应标注为修订视图，防止把事后计算结果误认为当时决策依据。

空间查询同样需要证据边界。淹没范围、影响人口或工程数量由特定模型和数据版本计算，界面点击统计数字时应能打开其空间范围、过滤条件和生成时间。截图只能作为沟通材料，正式结果应保留可查询数据和生成记录。

三维性能下降时优先保证业务信息可达。模型可降级LOD或切换二维地图，告警列表、数据表和处置入口仍应可用。通过这种渐进增强，三维场景不会成为整个业务系统的单点故障。

\subsection{实时同步与闭环控制}

一次完整闭环包括“感知—质检—状态更新—模型计算—风险解释—人工确认—指令下发—执行反馈—效果评估”。并非所有场景都允许自动控制：大坝安全和防洪调度通常需要权限、会商和人工确认。

\begin{figure}[htbp]
\centering
\begin{tikzpicture}[>=Stealth,node distance=.55cm]
\tikzset{dtstep/.style={draw=Primary,rounded corners=2pt,fill=PrimaryLight,
  minimum width=4.0cm,minimum height=.65cm,align=center,text=PrimaryDark}}
\node[dtstep] (sense) {1 感知与时间同步};
\node[dtstep,below=of sense] (qc) {2 质量控制与状态估计};
\node[dtstep,below=of qc] (model) {3 模型计算与不确定性评估};
\node[dtstep,below=of model] (decision) {4 风险解释与人工确认};
\node[dtstep,below=of decision] (act) {5 指令下发与执行回执};
\node[dtstep,below=of act] (verify) {6 效果校核与模型修正};
\draw[->,draw=Primary,thick] (sense)--(qc);
\draw[->,draw=Primary,thick] (qc)--(model);
\draw[->,draw=Primary,thick] (model)--(decision);
\draw[->,draw=Primary,thick] (decision)--(act);
\draw[->,draw=Primary,thick] (act)--(verify);
\draw[->,draw=Accent,thick] (verify.east) -- ++(2,0) |- (sense.east);
\end{tikzpicture}
\caption{数字孪生业务闭环}
\label{fig:dt-closed-loop}
\end{figure}

下面的短伪代码只表达状态更新主线。每个函数均代表可测试的服务接口，不虚构十余个未定义类。

\begin{lstlisting}[language=Python]
def update_twin(reading, repository, model, publisher):
    checked = validate_unit_time_quality(reading)
    repository.save_raw(checked)

    state = repository.load_state(checked.asset_id)
    state = merge_observation(state, checked)
    prediction = model.run(state)

    event = {
        "assetId": checked.asset_id,
        "stateVersion": state.version,
        "risk": prediction.risk,
        "modelVersion": model.version
    }
    repository.save_result(event)
    publisher.publish("twin.state.updated", event)
    return event
\end{lstlisting}

生产实现还需事务、幂等、超时、重试和审计。状态版本可防止旧结果覆盖新结果；模型版本和输入快照保证结果可复现。

\subsection{接口与事件契约}

数字孪生各层通过稳定契约协作。查询接口适合读取当前状态，事件适合通知“状态已经发生变化”，命令用于请求执行有业务副作用的动作。三者不能都包装成含混的“消息”。

状态更新事件至少包含事件标识、对象标识、发生时间、状态版本、模型版本和追踪标识：

\begin{lstlisting}
{
  "eventId": "01J5...",
  "eventType": "twin.state.updated.v1",
  "assetId": "DAM-001-BLOCK-07",
  "occurredAt": "2026-08-05T08:30:00+08:00",
  "stateVersion": 1842,
  "modelVersion": "seepage-2.3.1",
  "risk": "warning",
  "traceId": "5d0e..."
}
\end{lstlisting}

事件类型携带契约版本，新增可选字段保持向后兼容；删除或改变字段语义时发布新版本。消费者用\texttt{eventId}去重，用\texttt{stateVersion}拒绝过期更新，用\texttt{traceId}串联采集、模型和告警日志。

控制命令比状态事件要求更严格。命令包含发起人、审批记录、目标设备、期望状态、有效期和幂等键；执行端返回接收、拒绝、执行中、成功或失败状态。超时不等于未执行，调用方必须查询最终状态，避免重复启闭设备。

\begin{figure}[htbp]
\centering
\begin{tikzpicture}[>=Stealth,node distance=.9cm and 1.0cm]
\tikzset{svc/.style={draw=Primary,rounded corners=2pt,fill=PrimaryLight,
 minimum width=2.4cm,minimum height=.7cm,align=center,text=PrimaryDark}}
\node[svc] (gateway) {接入网关};
\node[svc,right=of gateway] (state) {状态服务};
\node[svc,right=of state] (model) {模型服务};
\node[svc,right=of model] (alert) {告警服务};
\node[svc,below=of model] (audit) {事件库与审计};
\draw[->,draw=Primary] (gateway)--node[above,font=\small]{观测事件}(state);
\draw[->,draw=Primary] (state)--node[above,font=\small]{状态版本}(model);
\draw[->,draw=Primary] (model)--node[above,font=\small]{模型结果}(alert);
\draw[->,draw=Accent] (gateway)--(audit);
\draw[->,draw=Accent] (state)--(audit);
\draw[->,draw=Accent] (model)--(audit);
\draw[->,draw=Accent] (alert)--(audit);
\end{tikzpicture}
\caption{状态更新的事件契约与审计链}
\label{fig:dt-event-contract}
\end{figure}

接口契约需通过自动化测试验证，包括必填字段、单位、时区、枚举、重复事件和乱序事件。契约文档与代码同步版本化，避免不同团队凭口头约定解释同一字段。

\subsection{大坝安全监测场景}

大坝场景可选渗压、位移、应变、温度、库水位和环境量作为观测。虚拟实体包含坝段几何、测点拓扑、统计基线和必要的机理模型。服务包括异常识别、趋势分析、空间联动和处置流程。

一次渗压异常不应只改变三维颜色。平台需展示原始值、质量码、历史趋势、相关库水位、模型残差、阈值依据和邻近测点状态，帮助值班人员判断传感器故障、环境影响或结构异常。确认后形成告警事件、处置工单和复核记录。

\begin{table}[htbp]
\centering
\caption{大坝数字孪生场景的输入、处理与输出}
\label{tab:dam-twin-scenario}
\begin{tabular}{p{3.1cm}p{4.5cm}p{5.2cm}}
\toprule
环节 & 内容 & 质量控制 \\
\midrule
输入 & 渗压、位移、库水位、温度 & 时间对齐、单位、缺测、突变检查 \\
状态更新 & 测点状态、坝段状态、关联工况 & 版本控制、空间拓扑校核 \\
模型服务 & 基线比较、回归/机理计算、风险规则 & 适用范围、残差、不确定性 \\
输出 & 风险等级、证据链、处置建议 & 人工确认、权限、审计 \\
反馈 & 工单结果、复测、模型修正 & 闭环时间、误报漏报复盘 \\
\bottomrule
\end{tabular}
\end{table}

\subsection{流域预报调度场景}

流域场景的空间范围更大，涉及降雨预报、河网汇流、水库调度和下游影响。平台应把“预报”和“调度”分开：预报模型给出未来流量及不确定性，调度模型比较约束条件下的方案，业务人员结合预案和实时信息作出决定。

方案比选至少展示目标、约束、模型版本、边界条件、关键断面过程和风险指标。任何调度建议都要保留输入快照；预报更新后重新计算时，新旧方案不能混用同一编号。

\subsection{云边协同与部署}

现场边缘节点适合完成协议接入、缓存、初步质检和断网续传；中心平台负责全局数据治理、模型调度和综合应用。控制指令采用最小权限、双向认证和可审计通道，不能因引入“数字孪生”而绕过原有安全规程。

\begin{figure}[htbp]
\centering
\begin{tikzpicture}[>=Stealth,node distance=.8cm and .9cm]
\tikzset{box/.style={draw=Primary,rounded corners=2pt,fill=PrimaryLight,
  minimum width=3.1cm,minimum height=.8cm,align=center,text=PrimaryDark}}
\node[box] (sensor) {现场设备\\传感器/闸门};
\node[box,right=of sensor] (edge) {边缘节点\\接入/质检/缓存};
\node[box,right=of edge] (bus) {中心接入\\API/消息总线};
\node[box,above right=of bus] (data) {数据与模型平台};
\node[box,below right=of bus] (app) {业务应用与会商};
\draw[<->,draw=Primary,thick] (sensor)--(edge);
\draw[<->,draw=Primary,thick] (edge)--(bus);
\draw[<->,draw=Primary,thick] (bus)--(data);
\draw[<->,draw=Primary,thick] (bus)--(app);
\draw[<->,draw=Accent,dashed,thick] (data)--(app);
\end{tikzpicture}
\caption{数字孪生水利平台云边协同部署}
\label{fig:dt-edge-cloud}
\end{figure}

部署设计需明确网络分区、离线时长、数据补传顺序、模型降级和灾备目标。例如模型服务不可用时，平台仍应展示经质检的实时数据和已批准阈值规则；不能把整个监控能力绑定在单一模型上。

\subsection{安全、权限与失效降级}

数字孪生汇集工程数据并可能连接执行设备，需把安全边界纳入架构。身份认证确认主体，授权限制其对象和操作范围；网络分区隔离办公、平台、现场控制等区域；传输和存储加密保护数据；审计记录关键查询、模型发布、告警确认和控制命令。

\begin{table}[htbp]
\centering
\caption{数字孪生平台的失效场景与降级策略}
\label{tab:dt-degradation}
\begin{tabular}{p{3.2cm}p{4.5cm}p{4.8cm}}
\toprule
失效场景 & 风险 & 预期降级 \\
\midrule
现场网络中断 & 实时状态停更 & 边缘缓存、标记数据龄期、恢复后按序补传 \\
关键测站缺测 & 模型输入不完整 & 降低可信度或停算，禁止静默填零 \\
模型服务超时 & 预报/分析不可用 & 保留实时展示和批准的静态规则 \\
消息重复或乱序 & 状态回退、重复告警 & 事件去重与状态版本检查 \\
三维模型加载失败 & 空间界面不可用 & 降级到二维地图、表格和告警列表 \\
身份服务故障 & 权限无法确认 & 高风险操作默认拒绝，保留只读应急入口 \\
\bottomrule
\end{tabular}
\end{table}

高风险控制采用职责分离：分析人员提出建议，授权人员审批，执行系统校验有效期和设备状态。系统不得因“自动化”跳过工程运行规程。紧急手动操作也要在恢复后补录原因和结果。

安全测试包括越权访问、令牌失效、重放、消息篡改、接口限流和审计完整性。三维前端同样属于攻击面，模型文件、纹理和属性数据要进行来源校验，避免把不可信脚本或敏感属性直接交付浏览器。

\subsection{可观测性、数据血缘与运行值守}

平台上线后，运行团队必须回答三个问题：当前服务是否可用，当前数据是否可信，当前模型结果能否用于业务。仅监控CPU和内存无法回答后两个问题，因此需要把服务、数据、模型和业务闭环统一纳入可观测体系。每条关键链路应配置服务级目标（Service Level Objective，SLO），并明确统计窗口、允许失败比例、告警责任人和处置时限。

\begin{table}[htbp]
\centering
\caption{数字孪生平台的分层可观测指标}
\label{tab:dt-observability}
\begin{tabular}{p{2.1cm}p{4.2cm}p{4.2cm}p{2.5cm}}
\toprule
层面 & 关键指标 & 典型异常 & 处置责任 \\
\midrule
服务 & 可用率、P95延迟、错误率、队列积压 & 接口超时、消费滞后、缓存击穿 & 平台运维 \\
数据 & 到达延迟、完整率、重复率、质量码分布 & 测站断报、时钟漂移、单位突变 & 数据值守 \\
模型 & 运行成功率、残差、漂移、置信区间覆盖 & 输入越界、参数失效、结果突跳 & 模型负责人 \\
业务 & 告警确认时长、工单闭环率、指令回执率 & 告警无人确认、工单长期挂起 & 业务值班 \\
安全 & 越权拒绝、异常登录、审计缺口 & 令牌滥用、批量探测、日志中断 & 安全管理员 \\
\bottomrule
\end{tabular}
\end{table}

一次渗压异常从传感器采集到处置完成，至少跨越测点、边缘网关、消息主题、清洗规则、状态估计、异常模型、告警服务、三维界面和工单系统。应为这条链路分配贯穿始终的事件标识或关联标识，使日志、指标和追踪记录可以按同一业务事件聚合。若只记录各系统内部编号，事故复盘时就难以证明某个告警使用了哪一批原始数据和哪个模型版本。

数据血缘记录数据集从来源到结果的变换关系。最小血缘节点包含来源标识、时间范围、处理规则版本、输入与输出校验摘要、执行状态和责任主体；模型结果还要关联模型卡、参数集和运行环境。血缘不是另存一份不可维护的说明文档，而应由采集、计算和发布流程自动写入元数据存储，并能从三维对象或告警详情反向查询。

\begin{figure}[htbp]
\centering
\begin{tikzpicture}[node distance=5mm and 7mm,
  item/.style={draw=Primary,rounded corners,fill=PrimaryLight,align=center,minimum width=2.25cm,minimum height=0.85cm,font=\small}]
\node[item] (raw) {原始测值\\批次与校验};
\node[item,right=of raw] (clean) {清洗结果\\规则版本};
\node[item,right=of clean] (state) {状态估计\\参数与质量};
\node[item,right=of state] (model) {模型结果\\模型卡版本};
\node[item,below=of model] (alarm) {告警事件\\阈值与证据};
\node[item,left=of alarm] (ticket) {处置工单\\人员与时限};
\node[item,left=of ticket] (review) {复盘归档\\结论与改进};
\draw[->,draw=Primary,thick] (raw)--(clean);
\draw[->,draw=Primary,thick] (clean)--(state);
\draw[->,draw=Primary,thick] (state)--(model);
\draw[->,draw=Primary,thick] (model)--(alarm);
\draw[->,draw=Primary,thick] (alarm)--(ticket);
\draw[->,draw=Primary,thick] (ticket)--(review);
\draw[->,draw=Accent,dashed,thick,bend left=25] (review) to node[above,font=\scriptsize]{改进规则/模型} (clean);
\end{tikzpicture}
\caption{从原始数据到处置复盘的血缘与证据链}
\label{fig:dt-lineage}
\end{figure}

值守界面应按影响范围和紧迫程度合并告警，避免同一根因产生大量重复提示。告警必须包含对象、发生时间、当前值、质量码、触发规则、推荐动作和确认入口；抑制、合并或关闭告警也要保留操作者与理由。交接班记录尚未恢复的数据源、降级中的模型、未闭环工单和临时权限，防止信息只停留在即时通信中。

当数据质量或模型表现持续低于阈值时，系统应自动降低结果可信等级，必要时退出自动推荐，但仍保留原始数据、二维地图和人工规程。恢复正常不能只看服务重新启动，还应完成积压补传、时间顺序校验、状态重算、告警去重和责任人确认。这样，可观测性才从“发现程序故障”扩展为“证明业务链路已经可信恢复”。

\subsection{验证、验收与持续改进}

数字孪生验收既包含软件质量，也包含数据和模型可信度。建议分为数据、空间、模型、服务、闭环和运维六类指标。

\begin{table}[htbp]
\centering
\caption{数字孪生平台验收检查项}
\label{tab:dt-acceptance}
\begin{tabular}{p{2.3cm}p{5.0cm}p{5.4cm}}
\toprule
类别 & 检查项 & 证据 \\
\midrule
数据 & 完整率、时效、单位、质量码、血缘 & 抽样记录与异常处置日志 \\
空间 & CRS、高程基准、控制点、模型对齐 & 校核报告与同名点误差 \\
模型 & 适用范围、参数、校准、回放结果 & 模型卡、版本和测试数据 \\
服务 & 可用性、延迟、错误处理、权限 & 压测、安全测试和监控 \\
闭环 & 告警到工单、指令到回执、责任追踪 & 演练记录与审计链 \\
运维 & 备份、降级、灾备、变更和回滚 & 运维手册与恢复演练 \\
\bottomrule
\end{tabular}
\end{table}

上线前使用历史典型洪水、设备故障和缺测场景回放；上线后持续统计模型残差、告警准确性、处置时长和数据质量。改进必须经过版本化、测试、审批和回滚准备。

\subsection{课程项目实施路线}

课程项目不要求一次实现完整流域数字孪生，可分四次迭代，每次形成可验收成果。第一轮建立对象编码、空间基准和数据样例；第二轮完成三维场景与状态绑定；第三轮接入一个模型或规则并形成事件；第四轮完成处置闭环、测试和演示。

\begin{table}[htbp]
\centering
\caption{数字孪生课程项目的四次迭代}
\label{tab:dt-course-iterations}
\begin{tabular}{p{1.8cm}p{4.0cm}p{4.1cm}p{3.3cm}}
\toprule
迭代 & 主要任务 & 交付物 & 验收重点 \\
\midrule
一 & 场景选择、对象编码、CRS与数据字典 & 范围说明、数据样例 & 边界和元数据完整 \\
二 & 加载GIS/三维模型、绑定对象状态 & 可交互页面、映射表 & 定位正确、异常可见 \\
三 & 实现质检、规则或简化模型 & 模型卡、结果事件 & 输入契约、版本可追踪 \\
四 & 告警处置、审计、测试和降级 & 演示系统、验收报告 & 闭环可复现、失败可降级 \\
\bottomrule
\end{tabular}
\end{table}

团队角色可分为数据/空间、前端场景、后端服务和测试集成，但接口与对象编码必须共同评审。每次迭代使用同一组场景用例，避免各成员分别制作无法拼接的演示。

最终演示应包含正常流、数据质量异常、模型不可用和无权限操作四类场景。验收报告记录输入文件、运行版本、操作步骤、预期和实际结果，使其他小组能够复现，而不仅是录制一段成功视频。

\subsection{发展重点}

值得持续关注的方向包括：模型可信度与不确定性表达、跨系统语义互操作、云边协同、面向事件的实时数据处理、人在回路的决策支持以及网络与数据安全。这些方向都能通过可测指标验证。本章不罗列“元宇宙”“区块链天然保证可信”等缺乏场景和证据的口号。

我国数字孪生水利建设强调流域防洪、水资源管理和工程运行等实际任务\cite{twinWater2023,twinProject2024}。平台规划应从可闭环的业务问题出发，逐步扩展数据、模型和服务，而不是先建设不可维护的全量高精度场景。

\section*{小结}

本章建立了从WebGL2坐标与渲染、Three.js模型加载、OGC服务、CGCS2000与高程基准，到倾斜摄影、BIM/IFC和数字孪生架构的完整链条。三维场景提供空间表达，数据底板保证时空一致，模型服务给出可追溯分析，业务闭环负责把结果转化为受控行动。只有四者共同工作，三维平台才具备持续支撑智慧水利业务的能力。

\section*{章末交付物}

提交一个水利工程三维场景方案：空间参考与高程基准说明；GIS、倾斜摄影和BIM模型清单；一段可运行的Three.js或Cesium加载代码；数字孪生五维映射、分层架构和一条业务闭环；以及数据、模型和系统验收表。

\section*{思考题与练习题}

\begin{enumerate}
\item[\textbf{客观题}]
\item MVP矩阵中的$V$表示（\quad）。A. 模型矩阵\quad B. 观察矩阵\quad C. 投影矩阵\quad D. 纹理矩阵
\item CGCS2000地理二维坐标系的EPSG代码是（\quad）。A. 4326\quad B. 3857\quad C. 4490\quad D. 4547
\item WMS以预定义瓦片矩阵为核心。（判断：对／错）
\item \texttt{Cesium3DTileset.fromUrl}应等待Promise完成后再加入场景。（判断：对／错）
\item 像控点的坐标来自外业测量，可参与空三约束。（判断：对／错）
\item 数字孪生数据维度只包含清洗算法，不包含实时和历史数据。（判断：对／错）

\item[\textbf{简答与设计题}]
\item 说明模型、世界、观察和裁剪坐标系的变换关系。
\item 比较WMS、WMTS和WFS的返回内容及水利应用场景。
\item 解释椭球高与1985国家高程基准正常高的区别，并列出数据入库所需元数据。
\item 说明倾斜摄影特征点、连接点、像控点和检查点的区别。
\item 结合IFC 4.3官方范围，说明闸门、坝体和廊道应如何避免错误实体映射。
\item 为某水库绘制数字孪生五维映射，说明模型结果如何进入人工确认和处置闭环。

\item[\textbf{实践题}]
\item 在一课时内完成一个GLB模型加载、定位和业务标识绑定，并记录包围盒与加载错误处理。
\item 使用5个以上检查点分别计算$m_x$、$m_y$、$m_h$，判断成果是否满足给定限差。
\item 以模板补全表\ref{tab:dt-acceptance}，为一个大坝渗压异常场景设计最小验收用例集。
\end{enumerate}
